\documentclass[12pt,a4paper]{article}

% ─── Packages ──────────────────────────────────────────────────────────────────
\usepackage[utf8]{inputenc}
\usepackage[T1]{fontenc}
\usepackage[turkish]{babel}
\usepackage{geometry}
\usepackage{graphicx}
\usepackage{hyperref}
\usepackage{listings}
\usepackage{xcolor}
\usepackage{tabularx}
\usepackage{booktabs}
\usepackage{enumitem}
\usepackage{fancyhdr}
\usepackage{titlesec}
\usepackage{mdframed}
\usepackage{amsmath}
\usepackage{float}

% ─── featurebox (tcolorbox yerine mdframed) ────────────────────────────────────
\newenvironment{featurebox}[1]{%
  \begin{mdframed}[
    backgroundcolor=bgweek,
    linecolor=borderpurple,
    linewidth=0.8pt,
    innerleftmargin=8pt,
    innerrightmargin=8pt,
    innertopmargin=6pt,
    innerbottommargin=6pt,
    roundcorner=4pt
  ]
  \raggedright
  \noindent\textbf{\small #1}\\[4pt]
}{%
  \end{mdframed}
}

% ─── notebox (tcolorbox yerine mdframed) ───────────────────────────────────────
\newenvironment{notebox}{%
  \begin{mdframed}[
    backgroundcolor=infoblue!8,
    linecolor=infoblue!60,
    linewidth=0.6pt,
    innerleftmargin=8pt,
    innerrightmargin=8pt,
    innertopmargin=5pt,
    innerbottommargin=5pt,
    roundcorner=4pt
  ]
}{%
  \end{mdframed}
}\geometry{margin=2.5cm}

% Küçük taşmaları önlemek için
\sloppy
\setlength{\emergencystretch}{3em}
\hbadness=10000
\hfuzz=2pt

% ─── Colors ────────────────────────────────────────────────────────────────────
\definecolor{fuchsia}{RGB}{217, 70, 239}
\definecolor{codeblue}{RGB}{100, 160, 255}
\definecolor{codegray}{RGB}{40, 40, 40}
\definecolor{lightgray}{RGB}{245, 245, 245}
\definecolor{darkbg}{RGB}{20, 20, 30}
\definecolor{successgreen}{RGB}{34, 197, 94}
\definecolor{warnyellow}{RGB}{234, 179, 8}
\definecolor{infoblue}{RGB}{59, 130, 246}
\definecolor{bgweek}{RGB}{248, 245, 255}
\definecolor{borderpurple}{RGB}{192, 132, 252}

% ─── Code Listing Style ────────────────────────────────────────────────────────
\lstset{
  basicstyle=\ttfamily\footnotesize,
  backgroundcolor=\color{lightgray},
  keywordstyle=\color{blue}\bfseries,
  commentstyle=\color{gray}\itshape,
  stringstyle=\color{red!60!black},
  numbers=left,
  numberstyle=\tiny\color{gray},
  numbersep=8pt,
  showstringspaces=false,
  breaklines=true,
  breakatwhitespace=false,
  breakindent=0pt,
  columns=flexible,
  postbreak=\mbox{\textcolor{gray}{$\hookrightarrow$}\space},
  frame=single,
  rulecolor=\color{gray!30},
  captionpos=b,
  tabsize=2,
}

% ─── Header / Footer ───────────────────────────────────────────────────────────
\pagestyle{fancy}
\fancyhf{}
\rhead{\textcolor{fuchsia}{\small CodeAlchemist — Hafta 5}}
\lhead{}
\rfoot{\small Sayfa \thepage}

% ─── Title Formatting ──────────────────────────────────────────────────────────
\titleformat{\section}[block]{\large\bfseries\color{fuchsia}}{}{0em}{}[\titlerule]
\titleformat{\subsection}[block]{\normalsize\bfseries\color{blue!70!black}}{}{0em}{}
\titleformat{\subsubsection}[block]{\small\bfseries\color{gray!70!black}}{}{0em}{}

% ─── Hyperlink Setup ───────────────────────────────────────────────────────────
\hypersetup{
  colorlinks=true,
  linkcolor=fuchsia,
  urlcolor=codeblue,
  citecolor=gray,
  pdftitle={CodeAlchemist 5. Hafta Değişiklikleri},
  pdfauthor={CodeAlchemist Geliştirme Ekibi},
}

% ───────────────────────────────────────────────────────────────────────────────
\begin{document}

% ─── Title Page ────────────────────────────────────────────────────────────────
\begin{titlepage}
  \centering
  \vspace*{2cm}
  {\Huge\bfseries\color{fuchsia} CodeAlchemist}\\[0.5em]
  {\Large\color{gray} Haftalık Geliştirme Raporu}\\[0.3em]
  {\large\itshape 5. Hafta — Uygulanan Değişiklikler}\\[2cm]

  \hrule height 1pt
  \vspace{0.5cm}
  \begin{tabular}{@{}ll@{}}
    \textcolor{successgreen}{\bfseries\large ✓} & \textbf{Özellik 1} — İnteraktif Onboarding Turu \\[4pt]
    \textcolor{successgreen}{\bfseries\large ✓} & \textbf{Özellik 2} — Mobil Deneyim İyileştirmeleri \\[4pt]
    \textcolor{successgreen}{\bfseries\large ✓} & \textbf{Özellik 3} — Geçmiş Arama \& Filtreleme \\[4pt]
    \textcolor{successgreen}{\bfseries\large ✓} & \textbf{Özellik 4} — AI Yanıt Kalitesi Geri Bildirimi \\[4pt]
    \textcolor{successgreen}{\bfseries\large ✓} & \textbf{Özellik 5} — Monaco Kod Editörü Entegrasyonu \\[4pt]
    \textcolor{successgreen}{\bfseries\large ✓} & \textbf{Özellik 6} — Kod Diff ve Uygula Sistemi \\[4pt]
    \textcolor{successgreen}{\bfseries\large ✓} & \textbf{Özellik 7} — Proje / Workspace Yönetimi \\
  \end{tabular}
  \vspace{0.5cm}
  \hrule height 1pt

  \vspace{2cm}
  \begin{tabular}{rl}
    \textbf{Hafta}   & 5. Hafta \\
    \textbf{Tarih}   & Mart 2026 \\
    \textbf{Ortam}   & React 19 + Flask 3 + SQLAlchemy + Monaco Editor \\
    \textbf{Durum}   & \textcolor{successgreen}{\bfseries Tamamlandı} \\
  \end{tabular}
\end{titlepage}

% ─── Table of Contents ─────────────────────────────────────────────────────────
\tableofcontents
\newpage

% ═══════════════════════════════════════════════════════════════════════════════
\section{5. Hafta Yapılan Değişiklikler — Genel Bakış}
% ═══════════════════════════════════════════════════════════════════════════════

Bu hafta toplamda \textbf{7 özellik modülü} uygulamaya alınmıştır. Değişikliklerin büyük bölümü frontend odaklıdır; Geri Bildirim (Özellik 4) ve Proje/Workspace (Özellik 7) backend tarafında yeni tablo ve endpoint gerektirmiştir.

\vspace{0.4cm}

{\footnotesize
\begin{tabularx}{\textwidth}{clXcc}
\toprule
\textbf{\#} & \textbf{Özellik} & \textbf{Etkilenen Dosyalar} & \textbf{FE} & \textbf{BE} \\
\midrule
1 & Onboarding Turu     & \texttt{OnboardingTour.jsx}, \texttt{App.jsx}, \texttt{index.css} & \textcolor{successgreen}{✓} & — \\
2 & Mobil UX            & \texttt{index.css}                                                & \textcolor{successgreen}{✓} & — \\
3 & Geçmiş Arama        & \texttt{HistoryList.jsx}                                          & \textcolor{successgreen}{✓} & — \\
4 & Geri Bildirim       & \texttt{ChatInterface.jsx}, \texttt{models.py}, \texttt{app.py}   & \textcolor{successgreen}{✓} & \textcolor{successgreen}{✓} \\
5 & Monaco Editör       & \texttt{MonacoCodeEditor.jsx}, \texttt{ChatInterface.jsx}         & \textcolor{successgreen}{✓} & — \\
6 & Kod Diff \& Apply   & \texttt{ChatInterface.jsx}                                        & \textcolor{successgreen}{✓} & — \\
7 & Proje/Workspace     & \texttt{ProjectManager.jsx}, \texttt{models.py}, \texttt{app.py} & \textcolor{successgreen}{✓} & \textcolor{successgreen}{✓} \\
\bottomrule
\end{tabularx}
}

\newpage

% ═══════════════════════════════════════════════════════════════════════════════
\section{Özellik 1: İnteraktif Onboarding Turu}
% ═══════════════════════════════════════════════════════════════════════════════

\begin{featurebox}{Etkilenen Dosyalar}
\textbf{Yeni:} \texttt{client/src/components/OnboardingTour.jsx}\\
\textbf{Güncellendi:} \texttt{client/src/App.jsx}, \texttt{client/src/index.css}
\end{featurebox}

\subsection{Gereksinim Analizi}
Yeni kullanıcılar uygulamaya ilk girdiklerinde boş ekranla karşılaşıyordu; ne yapacaklarını bilemiyorlardı. Kullanıcı elde tutma oranını artırmak ve ilk deneyimi iyileştirmek amacıyla 5 adımlı bir karşılama sihirbazı tasarlanmıştır.

\subsection{Ekran Görüntüleri}

\begin{figure}[H]
  \centering
  \resizebox{\linewidth}{!}{\includegraphics{onboarding_step1.png}}
  \caption*{\textbf{Adım 1} — Hoşgeldin ekranı}
\end{figure}

\begin{figure}[H]
  \centering
  \resizebox{\linewidth}{!}{\includegraphics{onboarding_step2.png}}
  \caption*{\textbf{Adım 2} — AI Model Seçimi}
\end{figure}

\subsection{Sihirbaz Adımları}

\begin{tabularx}{\textwidth}{clX}
\toprule
\textbf{Adım} & \textbf{Başlık} & \textbf{İçerik} \\
\midrule
1 & Hoşgeldin     & Uygulama tanıtımı ve genel bakış \\
2 & Model Seç     & AI model seçiminin nasıl yapılacağı \\
3 & Soru Sor      & Chat arayüzü ve dosya yükleme rehberi \\
4 & Kenar Çubuğu  & Geçmiş ve topluluk erişimi \\
5 & GitHub        & Repo bağlama entegrasyonu \\
\bottomrule
\end{tabularx}

\vspace{0.4cm}

UI elementleri şunlardır: ilerleme göstergesi (progress dots), \textit{Atla} butonu ve \textit{İleri →} butonu. Son adımda \textit{İleri} yerine \textit{Başla} gösterilir.

\subsection{Tamamlanma Takibi — localStorage}

Sihirbazın daha önce gösterilip gösterilmediği \texttt{localStorage} üzerinden takip edilir. Anahtar: \texttt{codealchemist\_onboarding\_done}.

\begin{lstlisting}[language=JavaScript, caption={OnboardingTour.jsx — gösterim ve tamamlanma kontrolü}]
const TOUR_KEY = 'codealchemist_onboarding_done';

// Sihirbaz gösterilmeli mi?
export const shouldShowOnboarding = () => !localStorage.getItem(TOUR_KEY);

// Kullanıcı tamamladığında işaretle
const handleComplete = () => {
  localStorage.setItem(TOUR_KEY, 'true');
  onComplete?.();
};
\end{lstlisting}

\subsection{App.jsx Entegrasyonu}

Sihirbaz iki durumda tetiklenir: ilk ziyarette ve başarılı kayıt (register) sonrasında.

\begin{lstlisting}[language=JavaScript, caption={App.jsx — onboarding state yönetimi}]
const [showOnboarding, setShowOnboarding] = useState(
  () => shouldShowOnboarding()
);

// Kayıt sonrası da tetikle
const handleAuthSuccess = (data) => {
  // ... auth işlemleri ...
  if (shouldShowOnboarding()) setShowOnboarding(true);
};

// JSX render:
{showOnboarding && (
  <OnboardingTour onComplete={() => setShowOnboarding(false)} />
)}
\end{lstlisting}

\subsection{Test Prosedürü}

Sihirbazı sıfırlamak ve yeniden test etmek için tarayıcı konsolunda aşağıdaki komut çalıştırılır, ardından sayfa yenilenir:

\begin{lstlisting}[language=bash, caption={Onboarding sıfırlama komutu}]
localStorage.removeItem('codealchemist_onboarding_done'); location.reload();
\end{lstlisting}

\newpage

% ═══════════════════════════════════════════════════════════════════════════════
\section{Özellik 2: Mobil Deneyim İyileştirmeleri}
% ═══════════════════════════════════════════════════════════════════════════════

\begin{featurebox}{Etkilenen Dosyalar}
\textbf{Güncellendi:} \texttt{client/src/index.css}
\end{featurebox}

\subsection{Gereksinim Analizi}
Chat arayüzü masaüstü için optimize edilmiş olmakla birlikte, 375px–768px arası ekran genişliklerinde aşağıdaki sorunlar tespit edilmiştir:

\begin{itemize}
  \item Mesaj balonları ekran dışına taşıyordu
  \item Mobil klavye açıldığında viewport kayıyordu
  \item Geri bildirim ve aksiyon butonlarının touch target alanı çok küçüktü
  \item Model karşılaştırma grid'i 2 sütunlu olduğundan taşıyordu
\end{itemize}

\subsection{Uygulanan CSS Değişiklikleri}

\begin{lstlisting}[language=CSS, caption={index.css — mobil medya sorguları}]
@media (max-width: 768px) {
  /* Mesaj baloncukları: daha geniş */
  .max-w-[85%] { max-width: 93% !important; }
  .max-w-[90%] { max-width: 95% !important; }

  /* Klavye açıkken viewport kaymasını engelle */
  .flex.flex-col.h-full { height: 100dvh; }

  /* Geri bildirim butonları — büyük touch target */
  .text-sm.px-2.py-1.rounded-lg {
    padding: 0.375rem 0.625rem !important;
    font-size: 1rem !important;
  }

  /* Model karşılaştırma grid — tek sütun */
  .grid.grid-cols-2 {
    grid-template-columns: 1fr !important;
  }
}

/* Onboarding için halka animasyonu */
@keyframes onboarding-pulse {
  0%, 100% { box-shadow: 0 0 0 0 rgba(217, 70, 239, 0.4); }
  50%       { box-shadow: 0 0 0 12px rgba(217, 70, 239, 0); }
}
\end{lstlisting}

\subsection{Değişiklik Özeti}

\begin{tabularx}{\textwidth}{lXl}
\toprule
\textbf{Sorun} & \textbf{Uygulanan Çözüm} & \textbf{Etkilenen Element} \\
\midrule
Mesaj taşması        & \texttt{max-width} 93–95\%'e çıkarıldı    & \texttt{.max-w-{[}85\%{]}}, \texttt{.max-w-{[}90\%{]}} \\
Viewport kayması     & \texttt{height: 100dvh} kullanıldı         & Chat container \\
Küçük butonlar       & Padding ve font büyütüldü                 & Feedback butonları \\
2-sütun grid taşması & \texttt{grid-template-columns: 1fr}        & Model Compare bileşeni \\
Onboarding efekti    & \texttt{@keyframes onboarding-pulse}       & Onboarding highlight \\
\bottomrule
\end{tabularx}

\begin{notebox}
\textbf{Not:} \texttt{100dvh} (dynamic viewport height), iOS Safari gibi tarayıcılarda klavye açıldığında küçülen gerçek viewport yüksekliğini temsil eder; \texttt{100vh} bu durumda statik kalır ve içerik kayar.
\end{notebox}

\newpage

% ═══════════════════════════════════════════════════════════════════════════════
\section{Özellik 3: Geçmiş Arama ve Filtreleme}
% ═══════════════════════════════════════════════════════════════════════════════

\begin{featurebox}{Etkilenen Dosyalar}
\textbf{Güncellendi:} \texttt{client/src/components/HistoryList.jsx}
\end{featurebox}

\subsection{Gereksinim Analizi}
Zamanla biriken konuşmalar arasında başlık bazında arama yapılamaması, kullanıcıların eski konuşmalarına dönmesini zorlaştırıyordu. Eklenen arama kutusu ile konuşmalar anlık olarak filtrelenebilmektedir.

\begin{figure}[H]
  \centering
  \resizebox{\linewidth}{!}{\includegraphics{history_search.png}}
  \caption*{Geçmiş arama kutusu ve tarih filtresi (Today / Week / Month)}
\end{figure}

\subsection{Mimari Karar: Tamamen Client-Side}

Filtreleme işlemi sunucuya hiçbir istek göndermeden, mevcut konuşma listesi üzerinde yapılır. Bu yaklaşımın avantajları:

\begin{itemize}
  \item \textbf{Sıfır gecikme} — her tuş vuruşunda anında sonuç
  \item \textbf{Ek API maliyeti yok} — backend'e yük binmez
  \item \textbf{Basit durum yönetimi} — yalnızca bir \texttt{useState} hook'u yeterlidir
\end{itemize}

\subsection{Filtreleme Mantığı}

\begin{lstlisting}[language=JavaScript, caption={HistoryList.jsx — client-side arama filtresi}]
const [searchText, setSearchText] = useState('');
const [filterPeriod, setFilterPeriod] = useState('all'); // 'all', 'today', 'week', 'month'

const filtered = (conversations || []).filter(item => {
  // Text search
  if (searchText.trim() && !(item.title || '').toLowerCase().includes(searchText.toLowerCase())) {
    return false;
  }

  // Date filter
  if (filterPeriod !== 'all') {
    const itemDate = new Date(item.created_at);
    const now = new Date();
    if (filterPeriod === 'today' && itemDate.toDateString() !== now.toDateString()) {
      return false;
    }
    if (filterPeriod === 'week') {
      const oneWeekAgo = new Date(now.setDate(now.getDate() - 7));
      if (itemDate < oneWeekAgo) return false;
    }
    if (filterPeriod === 'month') {
      const oneMonthAgo = new Date(now.setMonth(now.getMonth() - 1));
      if (itemDate < oneMonthAgo) return false;
    }
  }
  return true;
});
\end{lstlisting}

\subsection{Arama Kutusu Bileşeni}

\begin{lstlisting}[language=JavaScript, caption={HistoryList.jsx — arama input JSX}]
<div className="relative mb-3">
  <span className="absolute left-3 top-1/2 -translate-y-1/2 text-xs">
    ��
  </span>
  <input
    type="text"
    value={searchText}
    onChange={e => setSearchText(e.target.value)}
    placeholder="Search conversations..."
    className="w-full bg-gray-900/60 border border-gray-700 rounded-lg
               pl-8 pr-8 py-2 text-xs text-gray-200"
  />
  {searchText && (
    <button
      onClick={() => setSearchText('')}
      className="absolute right-2 top-1/2 -translate-y-1/2 text-gray-400"
    >
      ✕
    </button>
  )}
</div>

<div className="flex justify-around text-xs mb-3">
  {['all', 'today', 'week', 'month'].map(period => (
    <button
      key={period}
      onClick={() => setFilterPeriod(period)}
      className={`px-2 py-1 rounded-md ${filterPeriod === period ? 'bg-purple-600 text-white' : 'bg-gray-700 text-gray-300'}`}
    >
      {period.charAt(0).toUpperCase() + period.slice(1)}
    </button>
  ))}
</div>

{/* Sonuç bulunamazsa */}
{filtered.length === 0 && searchText && (
  <p className="text-xs text-gray-500 text-center py-4">
    No conversations matching "{searchText}"
  </p>
)}
\end{lstlisting}

\subsection{Özellik Özeti}

\begin{tabularx}{\textwidth}{lX}
\toprule
\textbf{Özellik} & \textbf{Davranış} \\
\midrule
Anlık filtreleme   & Her tuş vuruşunda listede anlık \texttt{includes()} kontrolü yapılır \\
Hızlı temizleme    & \textbf{✕} butonu ile arama sıfırlanır \\
Tarih Filtreleme   & \textit{Today, Week, Month} seçenekleriyle zamansal kısıtlama uygulanır \\
Boş durum mesajı   & \textit{No conversations matching "..."} gösterilir \\
\bottomrule
\end{tabularx}

\newpage

% ═══════════════════════════════════════════════════════════════════════════════
\section{Özellik 4: AI Yanıt Kalitesi Geri Bildirimi}
% ═══════════════════════════════════════════════════════════════════════════════

\begin{featurebox}{Etkilenen Dosyalar}
\textbf{Yeni (BE):} \texttt{server/models.py} → \texttt{Feedback} tablosu\\
\textbf{Yeni (BE):} \texttt{server/app.py} → 3 endpoint (\texttt{/api/feedback}, \texttt{/api/feedback/<id>}, \texttt{/api/feedback/detail})\\
\textbf{Güncellendi (FE):} \texttt{client/src/components/ChatInterface.jsx}
\end{featurebox}

\subsection{Gereksinim Analizi}
Kullanıcıların AI yanıtlarını değerlendirememesi, UX kalitesinin ölçülmesini engelliyordu. beğeni/beğenmeme sistemi hem anlık geri bildirim kanalı hem de ilerideki model ince ayarı (fine-tuning) için veri kümesi işlevi görmektedir.

\subsection{Veritabanı Modeli}

\begin{lstlisting}[language=Python, caption={models.py — Feedback tablosu}]
class Feedback(db.Model):
    """AI yanıtları için kullanıcı geri bildirimi (thumbs up/down)"""
    id         = db.Column(db.Integer, primary_key=True)
    history_id = db.Column(db.Integer, db.ForeignKey('history.id'))
    user_id    = db.Column(db.Integer,
                           db.ForeignKey('user.id'), nullable=True)
    rating     = db.Column(db.Integer, nullable=False)  # +1 veya -1
    created_at = db.Column(db.DateTime, default=datetime.utcnow)

    __table_args__ = (
        db.UniqueConstraint('history_id', 'user_id',
                            name='_feedback_unique_uc'),
    )
\end{lstlisting}

\textbf{Unique constraint} sayesinde her kullanıcı her mesaj için yalnızca bir oy verebilir. Aynı oya tıklamak oyu geri alır; farklı oya tıklamak değiştirir (toggle + güncelleme).

\subsection{API Endpoint'leri}

\begin{tabularx}{\textwidth}{llX}
\toprule
\textbf{Metod} & \textbf{Yol} & \textbf{Davranış} \\
\midrule
POST & \texttt{/api/feedback}            & Oy gönder / toggle / güncelle \\
GET  & \texttt{/api/feedback/<history\_id>} beğeni / beğenmeme sayısı ve kullanıcının mevcut oyu \\
\bottomrule
\end{tabularx}

\subsubsection{Toggle Mantığı (POST /api/feedback)}

\begin{lstlisting}[language=Python, caption={app.py — oy ekleme, geri alma ve güncelleme}]
if existing:
    if existing.rating == rating:
        db.session.delete(existing)    # Aynı oy -> geri al
    else:
        existing.rating = rating        # Farklı oy -> güncelle
else:
    db.session.add(Feedback(
        history_id=history_id,
        user_id=user_id,
        rating=rating
    ))

db.session.commit()
\end{lstlisting}

\subsubsection{Hata Detayı ve Kategori (POST /api/feedback/detail)}

Kullanıcı beğenmeme (dislike) verdiğinde, olumsuz oyun sebebini anlamak için bir detay formu açılır. Bu form aracılığıyla hata kategorisi ve opsiyonel yorum \texttt{FeedbackDetail} tablosuna kaydedilir.

\begin{lstlisting}[language=Python, caption={models.py — FeedbackDetail modeli}]
class FeedbackDetail(db.Model):
    id          = db.Column(db.Integer, primary_key=True)
    feedback_id = db.Column(db.Integer, db.ForeignKey('feedback.id'))
    category    = db.Column(db.String(100), nullable=False) # örn. 'Wrong Code'
    comment     = db.Column(db.Text, nullable=True)
\end{lstlisting}

\subsection{Frontend Uygulama}

Her AI yanıtının aksiyon çubuğunun sağına beğeni ve beğenmeme butonları eklenmiştir. Seçili oy yeşil/kırmızı ile vurgulanır.

\subsubsection{Optimistic Update Stratejisi}

Kullanıcı deneyimini iyileştirmek için oy, API yanıtı beklenmeden anında UI'a yansıtılır. API hatası durumunda önceki durum geri yüklenir.

\begin{lstlisting}[language=JavaScript, caption={ChatInterface.jsx — optimistic update ile feedback handler}]
const handleFeedback = async (historyId, rating) => {
  const previousRating = feedbacks[historyId];
  // Toggle: aynı oya tıklanırsa null yap
  const newRating = previousRating === rating ? null : rating;

  // 1. Anında UI güncelle (optimistic)
  setFeedbacks(prev => ({ ...prev, [historyId]: newRating }));

  try {
    await fetch(`${apiBase}/api/feedback`, {
      method: 'POST',
      headers: {
        'Content-Type': 'application/json',
        ...authHeaders,
      },
      body: JSON.stringify({ history_id: historyId, rating }),
    });
  } catch (err) {
    // 2. Hata durumunda önceki duruma dön
    setFeedbacks(prev => ({ ...prev, [historyId]: previousRating }));
  }
};
\end{lstlisting}

\subsubsection{Buton Durum Renkleri}

\begin{tabular}{lll}
\toprule
\textbf{Durum} & \textbf{Be\u011feni Rengi} & \textbf{Be\u011fenmeme Rengi} \\
\midrule
Oy verilmedi       & Nötr gri & Nötr gri \\
Beğeni seçili      & \textcolor{successgreen}{\textbf{Yeşil vurgulu}} & Nötr gri \\
Beğenmeme seçili   & Nötr gri & \textcolor{red}{\textbf{Kırmızı vurgulu}} \\
\bottomrule
\end{tabular}

\newpage

% ═══════════════════════════════════════════════════════════════════════════════
\section{Özellik 5: Monaco Kod Editörü Entegrasyonu}
% ═══════════════════════════════════════════════════════════════════════════════

\begin{featurebox}{Etkilenen Dosyalar}
\textbf{Yeni:} \texttt{client/src/components/MonacoCodeEditor.jsx}\\
\textbf{Güncellendi:} \texttt{client/src/components/ChatInterface.jsx}
\end{featurebox}

\subsection{Gereksinim Analizi}
Kullanıcı kodu yalnızca düz bir \texttt{<textarea>} alanına yapıştırabiliyordu. VS Code'un motorunu kullanan Monaco Editor entegrasyonu sözdizimi vurgulaması, otomatik tamamlama, satır numaraları, katlanabilir bloklar ve çoklu dil desteği sunmaktadır.

\begin{figure}[H]
  \centering
  \resizebox{\linewidth}{!}{\includegraphics{monaco_autocomplete.png}}
  \caption*{Monaco Editör — sözdizimi vurgulaması ve snippet tamamlama}
\end{figure}

\subsection{Bağımlılık}

\begin{lstlisting}[language=bash, caption={package.json — mevcut paket}]
"@monaco-editor/react": "^4.7.0"
\end{lstlisting}

\subsection{MonacoCodeEditor Bileşeni}

\begin{lstlisting}[language=JavaScript, caption={MonacoCodeEditor.jsx — temel yapı}]
import Editor from '@monaco-editor/react';

const MonacoCodeEditor = ({
  value, onChange, language = 'python',
  onLanguageChange, theme = 'dark', height = '200px'
}) => {
  return (
    <Editor
      height={height}
      language={language}
      value={value}
      onChange={(val) => onChange?.(val ?? '')}
      theme={theme === 'dark' ? 'vs-dark' : 'light'}
      options={{
        minimap: { enabled: false },
        fontSize: 13,
        fontFamily: "'JetBrains Mono', Consolas, monospace",
        fontLigatures: true,
        wordWrap: 'on',
        quickSuggestions: true,
        acceptSuggestionOnEnter: 'smart',
      }}
    />
  );
};
\end{lstlisting}

\subsection{Özel Tamamlama Sağlayıcısı (Completion Item Provider)}

Monaco Editor'ün açık kaynaklı API'si kullanılarak editöre \textbf{özel bir snippet tamamlama sağlayıcısı} entegre edilmiştir. Kullanıcı yazmaya başladığında, dile özgü anahtar kelimeler şablon (snippet) biçiminde önerilir. Python için tanımlanan temel snippet'ler: \texttt{print}, \texttt{def}, \texttt{class}, \texttt{for}, \texttt{if} ve \texttt{import}.

\begin{lstlisting}[language=JavaScript, caption={MonacoCodeEditor.jsx — özel completion provider kaydı}]
monaco.languages.registerCompletionItemProvider('python', {
  provideCompletionItems: (model, position) => {
    const suggestions = [
      {
        label: 'def',
        kind: monaco.languages.CompletionItemKind.Snippet,
        insertText: 'def ${1:func_name}(${2:args}):\n\t${3:pass}',
        insertTextRules:
          monaco.languages.CompletionItemInsertTextRule.InsertAsSnippet,
        documentation: 'Fonksiyon tanımı',
      },
      {
        label: 'class',
        kind: monaco.languages.CompletionItemKind.Snippet,
        insertText: 'class ${1:ClassName}:\n\tdef __init__(self):\n\t\t${2:pass}',
        insertTextRules:
          monaco.languages.CompletionItemInsertTextRule.InsertAsSnippet,
        documentation: 'Sınıf tanımı',
      },
      // print, for, if, import ... aynı yapıda eklenir
    ];
    return { suggestions };
  },
});
\end{lstlisting}

\begin{notebox}
\textbf{Not:} Tamamlama sağlayıcısı dil başına ayrı ayrı kaydedilir. SQL, JavaScript ve Python için bağımsız provider'lar tanımlanmış; her biri kendi anahtar kelime setini ve insert şablonunu içermektedir.
\end{notebox}

\subsection{ChatInterface.jsx Entegrasyonu}

Kullanıcı \textbf{Kod Edit\xf6r\xfc} butonuna tıklayarak Monaco editörüne, \textbf{⌨️ Plain} butonuna tıklayarak düz textarea'ya geçiş yapabilir:

\begin{lstlisting}[language=JavaScript, caption={ChatInterface.jsx — Monaco/Textarea toggle}]
const [useMonacoEditor, setUseMonacoEditor] = useState(false);

{useMonacoEditor ? (
  <MonacoCodeEditor
    value={question}
    onChange={setQuestion}
    language={monacoLanguage}
    onLanguageChange={setMonacoLanguage}
    theme="dark"
    height="180px"
  />
) : (
  <textarea value={question} onChange={...} />
)}
\end{lstlisting}

\subsection{Desteklenen Diller}

\begin{tabularx}{\textwidth}{XXXX}
\toprule
Python & JavaScript & TypeScript & Go \\
Rust   & Java       & C/C++      & SQL \\
HTML   & CSS        & JSON       & Markdown \\
\bottomrule
\end{tabularx}

\vspace{0.4cm}

\begin{notebox}
\textbf{Bundle Etkisi:} Monaco Editor lazy load ile yüklendiğinden başlangıç bundle boyutunu artırmaz. Editör DOM'a mount olduğunda dinamik yüklenir. Tahmini ek yük: ilk açılışta $\approx$2 MB.
\end{notebox}

\newpage

% ═══════════════════════════════════════════════════════════════════════════════
\section{Özellik 6: Kod Diff ve Uygula Sistemi}
% ═══════════════════════════════════════════════════════════════════════════════

\begin{featurebox}{Etkilenen Dosyalar}
\textbf{Güncellendi:} \texttt{client/src/components/ChatInterface.jsx}\\
\textbf{Bağımlılık:} \texttt{react-diff-viewer-continued@4.1.2} (mevcut)
\end{featurebox}

\subsection{Gereksinim Analizi}
AI bir kod düzeltmesi önerdiğinde, kullanıcı değişikliği elle kopyalamak zorundaydı. \textit{Apply Patch} butonu bu süreci tek bir tıklamaya indirerek iş akışını hızlandırır. Bu özellik, uygulamayı basit bir chat arayüzünden profesyonel bir \textbf{AI Kod Asistanına} (Cursor veya GitHub Copilot benzeri) dönüştüren kritik bir \textit{Quality of Life} (yaşam kalitesi) geliştirmesidir.

\clearpage
\begin{figure}[htbp]
\centering
\resizebox{\linewidth}{!}{\includegraphics{diff_apply.png}}
\caption*{Yan yana diff görünümü ve Apply Patch butonu}
\end{figure}

\subsection{Bağlam (Context) Yönetimi}

Bu özelliğin teknik olarak en büyük avantajı bağlam yönetimidir. CodeAlchemist'in çalışma prensibi gereği, editördeki kod her zaman bir sonraki sorunun bağlamı olarak sunucuya iletilir. \textit{Apply Patch} tıklandığında kod editörü anında güncellenir; böylece kullanıcı bir sonraki sorusunda AI'ya "bu kodu şu şekilde değiştir" dediğinde, AI zaten en son yamalanmış (güncel) kod üzerinden yanıt üretmeye devam eder.

\subsection{Diff Tetikleme Mekanizması}

Sistem, AI yanıtı içindeki özel marker'ları algılayarak diff görünümünü otomatik olarak tetikler. Yanıt \texttt{<<<OLD>>>} ve \texttt{<<<NEW>>>} etiketlerini içerdiğinde, bu iki blok arasındaki fark yan yana (split view) gösterilir.

\begin{lstlisting}[language=JavaScript, caption={ChatInterface.jsx — Marker bazlı diff tespiti}]
if (codeString.includes('<<<OLD>>>') && codeString.includes('<<<NEW>>>')) {
  const oldPart = codeString.split('<<<OLD>>>')[1].split('<<<NEW>>>')[0].trim();
  const newPart = codeString.split('<<<NEW>>>')[1].trim();
  // ... ReactDiffViewer render edilir
}
\end{lstlisting}

\subsection{Apply Patch İş Akışı}

\subsubsection{Adım 1 — Kod Bloklarını Çıkarma}

\begin{lstlisting}[language=JavaScript, caption={ChatInterface.jsx — kod bloklarını ayıklama}]
// Diff içermeyen standart kod bloklarını ayıklar
const extractCodeBlocks = (responseText) => {
  const regex = /```(\w+)?\n([\s\S]*?)```/g;
  // ...
};
\end{lstlisting}

\subsubsection{Adım 2 — Diff Görüntüleme}

\begin{lstlisting}[language=JavaScript, caption={ChatInterface.jsx — yan yana diff görünümü}]
import ReactDiffViewer from 'react-diff-viewer-continued';

<ReactDiffViewer
  oldValue={originalCode}   // Kullanıcının orijinal kodu
  newValue={suggestedCode}  // AI'nın önerdiği kod
  splitView={true}
  useDarkTheme={true}
  showDiffOnly={false}
/>
\end{lstlisting}

\subsubsection{Adım 3 — Patch Uygulama}

\textbf{Apply Patch} butonuna tıklandığında gerçekleşen işlemler:

\begin{enumerate}
  \item Kullanıcının aktif kod alanı (Monaco veya textarea) \texttt{setCode} fonksiyonu ile güncellenir.
  \item Değişiklik, ekranın üst kısmında beliren fuchsia renkli bir \textbf{toast bildirimi} ile görsel olarak onaylanır.
  \item Uygulanan kod, bir sonraki AI sorgusu için otomatik olarak bağlam (context) olarak ayarlanır.
\end{enumerate}

\newpage

% ═══════════════════════════════════════════════════════════════════════════════
\section{Özellik 7: Proje / Workspace Yönetimi}
% ═══════════════════════════════════════════════════════════════════════════════

\begin{featurebox}{Etkilenen Dosyalar}
\textbf{Yeni (FE):} \texttt{client/src/components/ProjectManager.jsx}\\
\textbf{Güncellendi (FE):} \texttt{client/src/App.jsx}, \texttt{client/src/components/ChatInterface.jsx}\\[2pt]
\textbf{Yeni (BE):} \texttt{server/models.py} → \texttt{Project}, \texttt{ProjectFile} tabloları\\
\textbf{Güncellendi (BE):} \texttt{server/app.py} → 8 endpoint + embedding tabanlı bağlam seçimi
\end{featurebox}

\subsection{Gereksinim Analizi}
Kullanıcılar tek bir konuşmada izole şekilde çalışabiliyordu. Gerçek dünya geliştirme senaryolarında birden fazla dosyanın (örn. \texttt{App.jsx} + \texttt{app.py}) aynı AI bağlamında sunulması gerekir. İlk sürümde tüm dosyalar doğrudan bağlama ekleniyordu; bu yaklaşım büyük projelerde gereksiz token tüketimine ve alakasız bağlam şişmesine neden oluyordu.

Bu haftaki güncelleme ile proje bağlamı \textbf{embedding tabanlı semantik seçim} ile güçlendirilmiştir. Sistem soru için en alakalı dosya parçalarını seçip \texttt{/api/ask} akışına ekler.

\begin{notebox}
\textbf{Süreç Filtresi:} Proje bağlamı enjeksiyonu doğrudan \texttt{/api/ask} pipeline'ı içinde yapılır. Uygulama, konuşmanın bir projeye bağlı olduğunu algıladığında ilgili dosyaları semantik olarak seçer ve AI'ya sistem mesajı olarak iletir.
\end{notebox}

\begin{notebox}
\textbf{Vurgu:} Proje sohbetinde semantik retrieval için öncelikli model \texttt{gemini-embedding-2-preview} olarak kullanılmaktadır. Model erişilemezse sistem \texttt{gemini-embedding-001} modeline düşer ve son çare olarak statik bağlam birleştirme yöntemine geri döner.
\end{notebox}

\subsection{Veritabanı Modelleri}

\begin{lstlisting}[language=Python, caption={models.py — Project ve ProjectFile tabloları}]
class Project(db.Model):
    id          = db.Column(db.Integer, primary_key=True)
    user_id     = db.Column(db.Integer, db.ForeignKey('user.id'))
    name        = db.Column(db.String(255), nullable=False)
    description = db.Column(db.Text, nullable=True)
    created_at  = db.Column(db.DateTime, default=datetime.utcnow)
    updated_at  = db.Column(db.DateTime, onupdate=datetime.utcnow)

    files = db.relationship('ProjectFile', backref='project',
                            lazy='dynamic', cascade='all, delete-orphan')


class ProjectFile(db.Model):
    id         = db.Column(db.Integer, primary_key=True)
    project_id = db.Column(db.Integer, db.ForeignKey('project.id'))
    name       = db.Column(db.String(255), nullable=False)
    content    = db.Column(db.Text, nullable=False, default='')
    language   = db.Column(db.String(50), default='plaintext')
    created_at = db.Column(db.DateTime, default=datetime.utcnow)
\end{lstlisting}

\texttt{cascade='all, delete-orphan'} sayesinde bir proje silindiğinde tüm dosyaları otomatik olarak temizlenir.

\subsection{API Endpoint'leri}

\begin{tabularx}{\textwidth}{llX}
\toprule
\textbf{Metod} & \textbf{Yol} & \textbf{Açıklama} \\
\midrule
GET    & \texttt{/api/projects}                  & Projeleri listele \\
POST   & \texttt{/api/projects}                  & Yeni proje oluştur \\
DELETE & \texttt{/api/projects/<id>}             & Projeyi sil \\
GET    & \texttt{/api/projects/<id>/files}       & Dosya listesi \\
POST   & \texttt{/api/projects/<id>/files}       & Dosya ekle \\
DELETE & \texttt{/api/projects/<id>/files/<fid>} & Dosyayı sil \\
GET    & \texttt{/api/projects/<id>/context}     & AI bağlamı (tüm dosyalar birleşik) \\
POST   & \texttt{/api/projects/<id>/semantic\_search} & Soruya göre semantik olarak en alakalı dosya parçalarını döndür (debug/inspection) \\
\bottomrule
\end{tabularx}

\subsection{AI Bağlamı Oluşturma (Embedding Destekli)}

\subsubsection{Aşama 1 — Chunking ve Embedding İndeksleme}

Proje dosyaları chunk'lara ayrılır, her chunk için embedding üretilir ve kısa süreli cache'te tutulur. Böylece aynı proje için tekrar eden sorularda yeniden embedding üretme maliyeti azaltılır.

\begin{lstlisting}[language=Python, caption={app.py — embedding tabanlı indeksleme (özet akış)}]
def _build_project_embedding_index(project):
  files = project.files.order_by(ProjectFile.name).all()
  signature = _project_signature(files)  # dosya değişimini izleme

  if cache_hit(signature):
    return cached_index

  chunks = []
  for pf in files:
    for chunk in _chunk_text(pf.content, chunk_size=1200, overlap=200):
      emb, model_used = _embed_text_with_fallback(
        chunk,
        task_type='RETRIEVAL_DOCUMENT'
      )
      if emb:
        chunks.append({
          'file': pf.name,
          'language': pf.language,
          'text': chunk,
          'embedding': emb,
          'model_used': model_used,
        })

  return persist_cache(signature, chunks)
\end{lstlisting}

\subsubsection{Aşama 2 — Query Embedding ve Top-K Seçim}

Kullanıcı sorusu \texttt{RETRIEVAL\_QUERY} task type ile embed edilir. Ardından cosine similarity ile en alakalı \texttt{top\_k} chunk seçilir.

\begin{lstlisting}[language=Python, caption={app.py — semantik hit çıkarımı (özet akış)}]
def get_project_semantic_hits(project, question, top_k=6):
  index_data = _build_project_embedding_index(project)
  query_emb, query_model = _embed_text_with_fallback(
    question,
    task_type='RETRIEVAL_QUERY'
  )

  scored = []
  for item in index_data['chunks']:
    score = _cosine_similarity(query_emb, item['embedding'])
    if score > 0:
      scored.append((score, item))

  scored.sort(key=lambda x: x[0], reverse=True)
  return scored[:top_k]
\end{lstlisting}

\subsubsection{Aşama 3 — \texttt{/api/ask} İçine Otomatik Enjeksiyon}

Konuşma bir projeye bağlıysa, semantik olarak seçilen chunk'lar sistem bağlamına eklenir. Embedding tarafında hata oluşursa mevcut statik yöntem (dosyaları birleştir) otomatik devreye girer.

\begin{lstlisting}[language=Python, caption={app.py — project context injection (özet akış)}]
if conversation.project_id:
  proj = db.session.get(Project, conversation.project_id)
  project_context = build_project_context_for_question(proj, question)

  if not project_context:
    project_context = build_static_project_context(proj)  # fallback

  github_context = project_context + "\n" + github_context
\end{lstlisting}

\subsection{Kullanıcı Deneyimi ve Doküman Desteği}

\subsubsection{Expand/Collapse Sohbet Listesi}
Sidebar üzerindeki projeler, o projeye ait konuşmaları hiyerarşik bir yapıda gösterir. Kullanıcı bir projeye tıkladığında o projenin konuşmaları expand/collapse mantığıyla listelenir.

\subsubsection{Bağımsız Model Seçici}
\texttt{ProjectWorkspace} içinde, global ayarlardan bağımsız bir model seçici bulunur. Bu sayede kullanıcı proje üzerinde çalışırken hızlıca model değiştirebilir.

\subsubsection{PDF ve DOCX Desteği (Base64 Pipeline)}
Binary dosyalar frontend tarafında base64 formatına çevrilerek sunucuya iletilir. Sunucu tarafında \texttt{pdfplumber} ve \texttt{python-docx} kütüphaneleri kullanılarak metin içeriği ayıklanır ve 20.000 karakter sınırı dahilinde projeye eklenir.

\subsection{Operasyonel Limitler}

\begin{tabularx}{\textwidth}{lXl}
\toprule
\textbf{Limit} & \textbf{Değer} & \textbf{Kaynak} \\
\midrule
Proje bağlamı kapasitesi & 12.000 karakter & \texttt{app.py:2341} \\
PDF/DOCX metin çıkarma sınırı & 20.000 karakter & \texttt{app.py:4993} \\
\bottomrule
\end{tabularx}

\subsection{Semantic Search Debug Endpoint'i}

Yeni endpoint, proje bazında semantik seçim sonucunu gözlemlemek için eklenmiştir:

\begin{itemize}
  \item \texttt{POST /api/projects/<id>/semantic\_search}
  \item Girdi: \texttt{query}, opsiyonel \texttt{top\_k}
  \item Çıktı: \texttt{query\_model}, \texttt{total\_chunks}, benzerlik skorları ve snippet listesi
\end{itemize}

Bu endpoint doğrudan ürün ekranına bağımlı değildir; operasyonel doğrulama, kalite analizi ve troubleshooting için tasarlanmıştır.

\begin{lstlisting}[language=Python, caption={app.py — /api/projects/<id>/context endpoint}]
@app.route('/api/projects/<int:project_id>/context')
@jwt_required()
def get_project_context(project_id):
  user_id = int(get_jwt_identity())
  project = Project.query.filter_by(
    id=project_id, user_id=user_id).first_or_404()
  files = project.files.order_by(ProjectFile.name).all()

  context_parts = [f"# Project: {project.name}"]
  for f in files:
      context_parts.append(
          f"\n## File: {f.name} ({f.language})"
          f"\n```{f.language}\n{f.content}\n```"
      )
  return jsonify({'context': '\n'.join(context_parts)})
\end{lstlisting}

\subsection{Frontend: ProjectManager Bileşeni}

\begin{figure}[H]
  \centering
  \resizebox{\linewidth}{!}{\includegraphics{project_screen.png}}
  \caption*{ProjectManager — proje oluşturma ekranı}
\end{figure}

\begin{lstlisting}[language=JavaScript, caption={ProjectManager.jsx — proje seçimi ve context enjeksiyonu}]
// Proje seçildiğinde App bileşenine bildirim gönder
const handleSelectAndClose = (project) => {
  onSelectProject?.(project);  // App.jsx'te activeProject güncellenir
  onClose?.();
};

// Not: Context prepend frontend'de değil backend /api/ask içinde yapılır.
// Backend artık embedding tabanlı alakalı chunk seçimi uygular.
// Frontend yalnızca project_id ile konuşma başlatır:
await fetch(`${apiBase}/api/conversations`, {
  method: 'POST',
  headers: { ...authHeaders, 'Content-Type': 'application/json' },
  body: JSON.stringify({ title, project_id: project.id })
});
\end{lstlisting}

\begin{lstlisting}[language=JavaScript, caption={App.jsx — Projects sekmesi entegrasyonu}]
{showProjectManager && (
  <ProjectManager
    apiBase={API_BASE}
    authHeaders={token ? { Authorization: `Bearer ${token}` } : {}}
    onSelectProject={(project) => setActiveProject(project)}
    activeProjectId={activeProject?.id}
    onClose={() => setShowProjectManager(false)}
  />
)}
\end{lstlisting}

\subsection{Frontend: Context Hits Kartı (ChatInterface)}

Proje sohbetlerinde kullanıcıya hangi bağlam parçalarının seçildiğini göstermek için sohbet balonunun altında \textbf{Context Hits} kartı eklenmiştir. Bu kart, görünürlük (explainability) katmanı sağlar.

\begin{figure}[H]
  \centering
  \resizebox{\linewidth}{!}{\includegraphics{project_screen2.png}}
  \caption*{Context Hits kartı — semantik olarak seçilen dosya parçaları ve benzerlik skorları}
\end{figure}

\begin{itemize}
  \item Yeni mesaj gönderildiğinde frontend, \texttt{/api/projects/<id>/semantic\_search} endpoint'ini çağırır
  \item Dönen \texttt{query\_model}, \texttt{total\_chunks} ve \texttt{hits} bilgileri ilgili chat turuna yazılır
  \item Kartta tekrar eden dosyalar tekilleştirilir; her dosya için en yüksek skorlu chunk gösterilir
\end{itemize}

\begin{lstlisting}[language=JavaScript, caption={ChatInterface.jsx — Context Hits kartı (özet akış)}]
const getUniqueContextHits = (hits = [], limit = 5) => {
  const bestByFile = new Map();
  for (const hit of hits) {
    const file = hit?.file || 'unknown';
    const score = Number(hit?.score || 0);
    const existing = bestByFile.get(file);
    if (!existing || score > Number(existing?.score || 0)) {
      bestByFile.set(file, hit);
    }
  }
  return Array.from(bestByFile.values())
    .sort((a, b) => Number(b?.score || 0) - Number(a?.score || 0))
    .slice(0, limit);
};

// render bölümünde
const uniqueHits = getUniqueContextHits(turn.semantic_context.hits, 5);
// dosya adı + chunk_index + skor gösterimi
\end{lstlisting}

\newpage

% ═══════════════════════════════════════════════════════════════════════════════
\section{5. Hafta Özet ve Değerlendirme}
% ═══════════════════════════════════════════════════════════════════════════════

\subsection{Değişiklik Matrisi}

\begin{tabularx}{\textwidth}{lXccc}
\toprule
\textbf{\#} & \textbf{Özellik} & \textbf{Yeni Dosya} & \textbf{FE} & \textbf{BE} \\
\midrule
1 & Onboarding Turu      & \texttt{OnboardingTour.jsx}               & \textcolor{successgreen}{✓} & — \\
2 & Mobil UX             & —                                         & \textcolor{successgreen}{✓} & — \\
3 & Geçmiş Arama         & —                                         & \textcolor{successgreen}{✓} & — \\
4 & Geri Bildirim (+/-)  & \texttt{Feedback} modeli                  & \textcolor{successgreen}{✓} & \textcolor{successgreen}{✓} \\
5 & Monaco Editör        & \texttt{MonacoCodeEditor.jsx}             & \textcolor{successgreen}{✓} & — \\
6 & Kod Diff \& Apply    & —                                         & \textcolor{successgreen}{✓} & — \\
7 & Proje/Workspace      & \texttt{Project}, \texttt{ProjectFile} modelleri & \textcolor{successgreen}{✓} & \textcolor{successgreen}{✓} \\
\bottomrule
\end{tabularx}

\subsection{Yeni Veritabanı Tabloları}

\begin{tabular}{lll}
\toprule
\textbf{Tablo} & \textbf{Amaç} & \textbf{Kısıt} \\
\midrule
feedback}      & beğeni/beğenmeme kayıtları    & UNIQUE (history\_id, user\_id) \\
\texttt{project}       & Workspace meta verisi & user\_id FK \\
\texttt{project\_file} & Dosya içerikleri      & project\_id FK, cascade delete \\
\bottomrule
\end{tabular}

\subsection{Yeni API Endpoint'leri}

5. haftada toplam \textbf{11 yeni endpoint} eklenmiştir:

\begin{tabularx}{\textwidth}{llX}
\toprule
\textbf{Grup} & \textbf{Endpoint Sayısı} & \textbf{İşlevler} \\
\midrule
Feedback API & 2 & POST (oy ver/toggle), GET (sayaç) \\
Feedback Detail API & 1 & POST (kategori + yorum) \\
Project API  & 8 & CRUD + dosya yönetimi + AI bağlamı + semantic search \\
\midrule
\textbf{Toplam} & \textbf{11} & \\
\bottomrule
\end{tabularx}

\subsection{Performans ve Bundle Etkisi}

Özellik 1, 2 ve 3: Saf CSS/JSX değişikliği, sıfır ek bağımlılık.\\
Özellik 4: Yalnızca \texttt{fetch} API'si, ek paket yok.\\
Özellik 5: Monaco lazy load ile yüklenir; başlangıç bundle'a etkisi yoktur. İlk açılışta $\approx$2 MB.\\
Özellik 7: Backend'de 2 yeni tablo, frontend'de 1 yeni bileşen.

\vspace{1cm}
\begin{center}
\textcolor{fuchsia}{\rule{10cm}{0.5pt}}\\[0.3em]
\textit{CodeAlchemist — Haftalık Geliştirme Raporu}\\
\textit{5. Hafta (7 Özellik) — Mart 2026}
\end{center}

\end{document}